\begin{Exercise}[label={ex:somma_multipli}]
    Si consideri una lista di interi. Scrivere una funzione che riceve come parametro la lista ed un intero M
    e che restituisce la somma dei soli valori della lista che sono multipli di M. Se la lista è vuota, la funzione
    deve restituire -1.
\end{Exercise}

\begin{Answer}[ref={ex:somma_multipli}]
  	\lstinputlisting[style=cstylesmall]{src/somma_multipli.c}
\end{Answer}